\begin{Application}[][10 min]
Soit $f$ la fonction numérique définie par : \quad
$f(x) = \frac{x^2 - 3}{2x^2 + 4}$

\begin{enumerate}
  \item Déterminer $D_f$ puis étudier la parité de la fonction $f$.
  \item Calculer les images des nombres suivants $0 \ ; 1 \ $ et $ -1$ par $f$.
  \item
  \begin{enumerate}
    \item Résoudre dans $\mathbb{R}$ l'équation : $f(x) = \frac{1}{2}$
    \item Montrer que, pour tout réel $x$ : $f(x) \leqslant \frac{1}{2}$
    \item Le nombre $\frac{1}{2}$ est-il un maximum de $f$ sur $\mathbb{R}$ ?
  \end{enumerate}
  \item
  \begin{enumerate}
    \item Montrer que pour tout $x$ de $\mathbb{R}$ : $f(x) \geqslant -\frac{3}{4}$
    \item Le nombre $-\frac{3}{4}$ est-il un minimum de $f$ sur $\mathbb{R}$ ?
  \end{enumerate}
\end{enumerate}
\end{Application}
